\documentclass[twocolumn,prl,superscriptaddress,nofootinbib]{revtex4-2}

\usepackage{amsmath,amssymb}
\usepackage{mathtools}
\usepackage{booktabs}
\usepackage{hyperref}

\newcommand{\R}{\mathbb{R}}
\newcommand{\C}{\mathbb{C}}
\newcommand{\Z}{\mathbb{Z}}
\newcommand{\Q}{\mathbb{Q}}
\newcommand{\ket}[1]{|#1\rangle}
\newcommand{\bra}[1]{\langle#1|}
\newcommand{\KS}{\mathrm{KS}}
\newcommand{\CK}{\mathrm{CK}\text{-}31}

\usepackage{amsthm}
\newtheorem{conjecture}{Conjecture}

\begin{document}

\title{Graph isomorphism of 31-vertex Kochen--Specker sets\\across coordinate alphabets in dimension three}

\author{Michael Kernaghan}
\affiliation{Pacific Quantum Systems, Vancouver, Canada}

\date{March 2026}

\begin{abstract}
We present a computational survey of coordinate alphabets for Kochen--Specker (KS) constructions in dimension~3, testing quadratic fields, cyclotomic fields, and cubic extensions. KS-uncolorable sets were found only when the coordinate alphabet (after closure) contains elements enabling \emph{modulus-2 cancellations} ($|x|^2 = 2$) or \emph{phase cancellations} ($1 + \omega + \omega^2 = 0$). Six discrete algebraic islands are identified. In three alphabet-based search instances that yielded 31-ray KS subsets---integer, rational, and mixed-alphabet origins---the minimized 31-ray subsets are mutually graph-isomorphic, verified by the VF2 algorithm. We note that the rational and mixed alphabets contain the integer alphabet, so this observation may reflect rediscovery of an embedded Conway--Kochen ($\CK$) configuration rather than independent convergence. We conjecture that $\CK$ is the unique minimal KS set in dimension~3.
\end{abstract}

\maketitle

\section{Introduction}

The smallest known Kochen--Specker (KS) set in dimension~3 is the 31-vector Conway--Kochen set~\cite{ConwayKochen,Peres1993}, constructed from integer coordinates.  (We use the ``original'' KS definition~\cite{LiBrightGanesh2024}: every ray belongs to at least one orthogonal triple, no triple admits a $\{0,1\}$-coloring consistent with unit sum.)  The Peres construction~\cite{Peres1991} uses $\Z[\sqrt{2}]$ coordinates (33~vectors), and Cabello's recent simplest KS set~\cite{Cabello2025simplest} uses the Eisenstein integers $\Z[\omega]$ (33~vectors).  Gould and Aravind~\cite{GouldAravind2010} showed that the Peres and Penrose 33-vector sets---despite different coordinate fields---share the same orthogonality graph, and parametrized a continuous family with the constraint $|c|^2 = 2$.  The best proven lower bound is 24~vectors~\cite{LiBrightGanesh2024,KirchwegarPeitlSzeider2023}. A natural question arises: what algebraic property determines whether a number field supports KS sets, and is the 31-vector minimum specific to the integer field or a universal barrier?

We present a computational survey across coordinate alphabets drawn from algebraic number rings, showing that modulus-2 cancellation identities correlate with KS-uncolorability, and that three alphabet-based search instances yielding 31-ray KS subsets all produce the same graph.

\section{The modulus-2 cancellation boundary}

For each algebraic number ring, we form a \emph{coordinate alphabet} $\mathcal{A} = \{0, \pm 1, \pm\alpha, \pm\bar\alpha, \ldots\}$---a finite set of values from which ray coordinates are drawn (called ``vector components'' by Pavicic et al.~\cite{Pavicic2019}).  The relevant object is the alphabet, not the ambient field: the same field can host different alphabets with different KS outcomes (e.g.\ $\Z[\sqrt{-3}]$ and $\Z[\omega]$ both lie in $\Q(\sqrt{-3})$, but only the Eisenstein alphabet $\Z[\omega]$ supports KS sets in our search).  Two rays are orthogonal when $\sum_i \bar{u}_i v_i = 0$; for coordinates from a finite alphabet, this requires a \emph{cancellation identity}---an algebraic zero-sum that makes the dot product vanish exactly.  For the integer alphabet the key identity is $2 - 1 - 1 = 0$; for $\Z[\sqrt{2}]$ it is $\alpha\bar\alpha - 1 - 1 = 0$ where $|\alpha|^2 = 2$; for the Eisenstein integers, $1+\omega+\omega^2=0$.  Each identity enables exact orthogonal triples (bases), and when enough bases interlock, no consistent $\{0,1\}$-coloring exists.  (The trivial cancellation $1+(-1)=0$ in the minimal alphabet $\{0,\pm 1\}$ generates only 13 projective rays---far too few for a KS set.)  Our approach is to systematically construct alphabets from algebraic number rings, identify their cancellation identities, and test whether the resulting ray sets are KS-uncolorable.

We generate all rays in $\C P^2$ with coordinates in $\mathcal{A}$, close under Hermitian orthogonal completion (iterating the cofactor formula $w_k = (-1)^{k+1}\det[\bar{u}_{k+1},\bar{u}_{k+2};\bar{v}_{k+1},\bar{v}_{k+2}]$ until the ray pool stabilizes), and test KS-uncolorability via SAT (Glucose4 solver~\cite{AudemardSimon2018}, PySAT). Minimization uses randomized greedy deletion (500 trials, seed 42) and does not guarantee minimality; reported sizes are upper bounds.

Table~\ref{tab:boundary} summarizes results across 40+ alphabets drawn from different rings. The empirical pattern: KS-uncolorable sets were found only when the coordinate alphabet (after closure) contains elements enabling \emph{modulus-2 cancellations} ($|x|^2 = 2$) or \emph{phase cancellations} ($1 + \omega + \omega^2 = 0$).

\begin{table}[t]
\centering
\caption{Representative results across coordinate alphabets. The column $|\alpha|^2$ shows the Hermitian squared modulus of the ring generator. KS sets were found only when the effective coordinate set (after completion) contains modulus-2-type cancellations. ``Min'' is the smallest KS subset found (upper bound on true minimum).}
\label{tab:boundary}
\begin{tabular}{lccc}
\toprule
Coordinate ring & $|\alpha|^2$ & KS? & Min \\
\midrule
$\Z$ (integers, $\alpha=2$) & -- & Yes & 31 \\
$\Z[\sqrt{2}]$ & 2.000 & Yes & 33 \\
$\Z[\omega]$ (Eisenstein) & 1.000 & Yes & 33 \\
$\Z[\sqrt{-2}]$ & 2.000 & Yes & 33 \\
$\Z[(1{+}\sqrt{-7})/2]$ & 2.000 & Yes & 43 \\
$\Q(\varphi)$ (golden)$^\dagger$ & 2.618 & Yes & 52 \\
\midrule
$\Z[\sqrt{3}]$ & 3.000 & No & -- \\
$\Z[\sqrt{5}]$ & 5.000 & No & -- \\
$\Z[\sqrt{-3}]$ ($\neq \Z[\omega]$) & 3.000 & No & -- \\
$\Z[\sqrt{-5}]$ & 5.000 & No & -- \\
$\Q(\sqrt[3]{2})$ basic & 1.587 & No$^\ddagger$ & -- \\
$\Q(\sqrt[3]{2})$ extended & 2.520 & Yes & 60 \\
\bottomrule
\end{tabular}
\begin{flushleft}
\footnotesize $^\dagger$Cross-product completion generates coordinates with $|x|^2 \leq 2$ (e.g.\ $1/\varphi$ satisfies $|1/\varphi|^2 < 1$).\\
$^\ddagger$Basic alphabet colorable; extended alphabet $\{0,\pm 1,\pm\sqrt[3]{2},\pm\sqrt[3]{4}\}$ is uncolorable via completion.
\end{flushleft}
\end{table}

The mechanism is cancellation identities: algebraic zero-sums among alphabet elements that make dot products vanish exactly.  Across all tested alphabets, KS-uncolorability was found only when one of two types of cancellation is available: (A)~\emph{modulus-2 cancellation}---a generator $x$ with Hermitian squared modulus $|x|^2 = 2$, enabling zero-sums such as $\alpha\bar\alpha - 1 - 1 = 0$ (integers, Peres, Heegner-7); or (B)~\emph{phase cancellation}---a root-of-unity sum vanishing exactly, as in $1+\omega+\omega^2 = 0$ (Eisenstein).  These are qualitatively different: modulus-2 cancellations balance two small terms against one larger term, while phase cancellations sum three equal-magnitude complex terms to zero.  For alphabets where neither mechanism is available within our bounded search, we found orthogonal pairs but the resulting hypergraphs were KS-colorable. (Note: orthogonal triples do exist---e.g., $\Z[\sqrt{3}]$ produces 15 bases from 91 rays---but the hypergraph is colorable.)

The golden ratio alphabet illustrates an important subtlety: $|\varphi|^2 \approx 2.618 > 2$, but completion generates the reciprocal $1/\varphi$ with $|1/\varphi|^2 < 1$, introducing modulus-2-type cancellations into the effective coordinate set. The relevant criterion is not the generator's squared modulus alone but whether the \emph{generated coordinate set} (after closure) admits sufficient cancellations.

Crucially, the obstruction for $|\alpha|^2 \geq 3$ alphabets is \emph{not} the absence of orthogonal triples. The alphabets from $\Z[\sqrt{3}]$, $\Z[\sqrt{5}]$, and $\Z[\sqrt{-5}]$ all produce substantial pools (133 rays, 94 bases each after completion), yet remain KS-colorable. The bases exist but are not configured to force uncolorability---the cancellation structure is insufficient.

Additional evidence: among cyclotomic fields $\Q(e^{2\pi i/n})$ for $n \leq 30$, KS-uncolorability was observed when $6 \mid n$---precisely the values for which the cyclotomic ring contains the Eisenstein integers $\Z[\omega]$ as a subring (since $\omega = e^{2\pi i/3}$ and $-1 = e^{i\pi}$ both lie in $\Z[e^{2\pi i/n}]$ when $6 \mid n$).

\section{Six algebraic islands}

Among all tested fields, six algebraic islands were found to support KS sets:

\begin{table}[t]
\centering
\caption{The six algebraic KS islands found in dimension~3. ``Min'' values are upper bounds from greedy deletion. Graph type indicates isomorphism class of the smallest found set.}
\label{tab:islands}
\begin{tabular}{lccl}
\toprule
Island & Min & Bases & Graph type \\
\midrule
Integer ($\CK$) & 31 & 17 & Unique (31v) \\
Peres & 33 & 16 & Norm-2 (33v) \\
Eisenstein & 33 & 14 & Eisenstein (33v) \\
$\Z[\sqrt{-2}]$ & 33 & 16 & $\cong$ Peres \\
Heegner-7 & 43 & 23 & New \\
Golden ($\Q(\varphi)$) & 52 & -- & New \\
\bottomrule
\end{tabular}
\end{table}

The Peres and $\Z[\sqrt{-2}]$ islands produce graph-isomorphic minimal sets (verified by VF2), yielding five distinct graph types. Non-isomorphism among the remaining pairs is established by differing degree sequences. The Heegner-7 and golden ratio islands are new constructions not found in the four previously known algebraic KS families; explicit ray coordinates and construction details are provided as Supplemental Material~\cite{SupplementalMaterial}.

\section{Graph universality of $\CK$}

In three alphabet-based search instances that yielded 31-ray KS subsets, the minimized subsets are mutually graph-isomorphic.

The three search instances use different starting alphabets:
\begin{enumerate}
\item Integer alphabet $\{0, \pm 1, \pm 2\}$ (109 rays after completion)
\item Rational $\Z[\tfrac{1}{2}]$ alphabet $\{0, \pm 1, \pm 2, \pm\tfrac{1}{2}\}$ (109 rays)
\item Mixed modulus-2: $\{0, \pm 1, \pm 2, \pm\sqrt{2}, \pm\omega, \pm\bar\omega\}$ (489 rays after completion)
\end{enumerate}

\emph{Important caveat}: the rational alphabet is projectively equivalent to the integer alphabet (verified by canonicalizing all rays to have first nonzero coordinate equal to~1: both alphabets produce the same set of 109 projective rays), and the mixed alphabet contains the integer alphabet as a subset. We cannot rule out that the minimization procedure simply rediscovers the embedded CK configuration in each case. This observation is therefore consistent with---but does not prove---a uniqueness claim.

All three resulting smallest-found sets have identical invariants:

\begin{table}[h]
\centering
\caption{Graph and hypergraph invariants of all 31-vector KS sets found. All entries are identical across all three constructions.}
\label{tab:universality}
\begin{tabular}{lc}
\toprule
Invariant & Value \\
\midrule
Vertices & 31 \\
Edges (orthogonal pairs) & 71 \\
Bases (orthogonal triples) & 17 \\
Degree sequence & $(3^4, 4^{14}, 5^8, 6^3, 8^2)$ \\
Triangles & 17 \\
$\mathrm{tr}(A^4)$ & 1250 \\
Spectral radius & 4.9443 \\
WL-1 test & Indistinguishable \\
\bottomrule
\end{tabular}
\end{table}

We verified pairwise graph isomorphism for all $\binom{3}{2} = 3$ pairs using the VF2 algorithm~\cite{CordellaEtAl2004} (via NetworkX), which provides a definitive yes/no answer and an explicit vertex mapping. All three pairs are isomorphic. In dimension~3, every triangle in the orthogonality graph is already a basis triple (any two orthogonal rays in $\C^3$ determine a unique third), so graph isomorphism implies basis hypergraph isomorphism. The same orthogonality structure thus emerges from all three search instances, though as noted above, the rational and mixed searches may simply be rediscovering the embedded integer CK configuration. Explicit ray coordinates, edge lists, basis lists, and VF2 certificates are provided as Supplemental Material~\cite{SupplementalMaterial}.

\section{Discussion}

The graph isomorphism of the three 31-ray subsets has two potential implications, subject to the caveat that the three searches may not be genuinely independent.

Note that at 33~vectors, multiple non-isomorphic KS graphs coexist: the Sch\"utte set~\cite{Schutte1965,GouldAravind2010} has a different orthogonality graph from the Peres/Penrose family~\cite{GouldAravind2010,Cabello2025simplest}. That the 31-ray subsets found by our minimization converge to a single graph type is therefore at least noteworthy, even if the evidence for true universality remains limited.

First, this is consistent with 31 being the true minimum for KS sets in dimension~3: no tested alphabet breaks the barrier. However, our minimization is heuristic (500 greedy trials), not exhaustive, so we cannot exclude the possibility of other 31-ray KS graphs within the same parent configurations.

\begin{conjecture}[Optimality]
No KS set with 30 or fewer rays exists in dimension~3.
\end{conjecture}

\begin{conjecture}[Uniqueness]
$\CK$ is the unique 31-ray KS graph up to isomorphism.
\end{conjecture}

\noindent These are logically independent claims: optimality could hold with multiple non-isomorphic 31-ray KS graphs, and uniqueness could hold even if a 30-ray KS set existed in a different construction family.

The proven lower bound is 24~\cite{LiBrightGanesh2024,KirchwegarPeitlSzeider2023}; closing the gap between 24 and 31 remains the central open problem.  These conjectures are supported by: the rigidity result of Trandafir and Cabello~\cite{TrandafirCabello2025}, who showed that every bipartite perfect quantum strategy based on 31 rays yields a unique strategy up to local isometry (a result about quantum strategies, not directly about graph combinatorics, but suggestive of structural uniqueness); our verification that 31 is the minimum found within the integer island across all alphabet sizes tested; the VF2-verified graph isomorphism across three alphabet-based search instances at both graph and hypergraph level; and the failure of 30{,}000+ random KS-uncolorable 3-uniform hypergraphs on 31 vertices (generated by randomized greedy construction and tested for realizability in $\R^3$ via SAT-encoded orthogonality constraints; see Supplemental Material~\cite{SupplementalMaterial}) to achieve geometric realizability.  Notably, the algebraic alphabet method falls outside both construction families analyzed in Trandafir and Cabello's Appendix~B (dimension-lifting and concatenation), which they showed cannot produce rigid KS sets in~$\C^3$.  The algebraic approach, constrained by number-theoretic cancellation identities rather than geometric concatenation, circumvents this obstruction and yields rigid sets (unique up to unitary equivalence in~$\C^3$, in the sense of Trandafir and Cabello~\cite{TrandafirCabello2025}) in four of six islands tested.

Second, the modulus-2 cancellation pattern provides a \emph{predictive heuristic}: given a coordinate alphabet, one can estimate whether KS sets are likely by checking whether the alphabet supports modulus-2-type or phase-type cancellation identities.  An alphabet must carry \emph{some} cancellation identity to produce orthogonal triples at all; the empirical observation is that only identities of modulus-2 or phase type appear to generate enough interlocking bases for KS-uncolorability within our bounded search.  While this does not constitute a proof (the search space is not exhausted), the heuristic has correctly predicted the outcome for every alphabet tested, providing a useful filter for future searches.

For the restricted family of two-element alphabets $\{0, \pm 1, \pm x\}$, enumerating all three-term zero-sums from the product set $\{0, \pm 1, \pm x, \pm |x|^2\}$ yields exactly four non-trivial identities---$x = 2$, $|x|^2 = 2$, $|x|^2 = x + 1$, and $x^2 + x + 1 = 0$---each corresponding to a known island.  This enumeration is complete for two-element alphabets.  Our coverage extends to quadratic fields ($\Q(\sqrt{d})$ for $d = 2, \ldots, 30$; all nine class-1 imaginary quadratic fields) and cyclotomic fields ($n \leq 30$).  The principal open question is whether multi-element alphabets from higher-degree fields harbor cancellation identities not reducible to modulus-2 or phase type.  The cubic field $\Q(\sqrt[3]{2})$ already produces a KS set at 60~vectors, suggesting that the pattern may extend but with rapidly increasing cost; quartic and higher extensions remain unexplored.

\textit{Verification methodology.}---This paper was developed using an LLM-assisted workflow~\cite{Schwartz2026}.  All computational claims (VF2 isomorphism, SAT-based uncolorability, negative controls) are reproducible via scripts with deterministic seeding; key results were cross-verified using adversarial LLM peer review.

\begin{acknowledgments}
The author thanks Ad\'an Cabello for discussions on the algebraic structure of KS sets.  Claude (Anthropic) performed all computational work including VF2 isomorphism testing, SAT verification, and manuscript preparation; cross-verification was performed using GPT (OpenAI).  The author directed the research, validated results, and accepts full responsibility for the scientific content and integrity of this paper.
\end{acknowledgments}

\begin{thebibliography}{15}

\bibitem{KochenSpecker1967}
S.~Kochen and E.~P.~Specker,
J.\ Math.\ Mech.\ \textbf{17}, 59 (1967).

\bibitem{ConwayKochen}
J.~H.~Conway and S.~Kochen,
reported in A.~Peres, \textit{Quantum Theory: Concepts and Methods} (Kluwer, 1993).

\bibitem{Peres1991}
A.~Peres,
J.\ Phys.\ A \textbf{24}, L175 (1991).

\bibitem{Peres1993}
A.~Peres,
\textit{Quantum Theory: Concepts and Methods} (Kluwer, 1993).

\bibitem{Cabello2025simplest}
A.~Cabello,
Phys.\ Rev.\ Lett.\ (2025); arXiv:2508.07335.

\bibitem{CSW2014}
A.~Cabello, S.~Severini, and A.~Winter,
Phys.\ Rev.\ Lett.\ \textbf{112}, 040401 (2014).

\bibitem{TrandafirCabello2025}
S.~Trandafir and A.~Cabello,
Phys.\ Rev.\ A \textbf{111}, 052204 (2025); arXiv:2501.11640.

\bibitem{AudemardSimon2018}
G.~Audemard and L.~Simon,
in \textit{Proceedings of IJCAI-18}, 5093 (2018).

\bibitem{CordellaEtAl2004}
L.~P.~Cordella, P.~Foggia, C.~Sansone, and M.~Vento,
IEEE Trans.\ Pattern Anal.\ Mach.\ Intell.\ \textbf{26}, 1367 (2004).

\bibitem{Pavicic2019}
M.~Pavicic, N.~D.~Megill, B.~C.~Waegell, and P.~K.~Aravind,
Sci.\ Rep.\ \textbf{9}, 6765 (2019).

\bibitem{SupplementalMaterial}
See Supplemental Material for explicit ray coordinates, edge lists, basis lists, VF2 isomorphism certificates, SAT solver details, and negative control data for all constructions reported here.

\bibitem{Schutte1965}
K.~Sch\"utte,
reported in J.~Bub, \textit{Interpreting the Quantum World} (Cambridge University Press, 1997), p.~133.

\bibitem{GouldAravind2010}
P.~K.~Gould and P.~K.~Aravind,
Found.\ Phys.\ \textbf{40}, 1096 (2010).

\bibitem{LiBrightGanesh2024}
Z.~Li, C.~Bright, and V.~Ganesh,
in \textit{Proceedings of IJCAI-24} (2024); arXiv:2306.13319.

\bibitem{KirchwegarPeitlSzeider2023}
M.~Kirchweger, A.~Peitl, and S.~Szeider,
in \textit{Proceedings of AAAI-23}, 4119 (2023).

\bibitem{Schwartz2026}
M.~D.~Schwartz,
``Vibe physics: the AI grad student,''
Anthropic Research (2026).

\end{thebibliography}

\end{document}
